\section{Preliminaries}

\paragraph{From Multi-Step to One-Step Action Generation.}
In conditional imitation learning, given a dataset of expert demonstrations \(\mathcal{D} = \{(o_i, a^*_i)\}_{i=1}^N\), 
the objective is to learn a policy mapping an observation \(o \in \mathcal{O}\) to a expert action chunk \(a^*\) 
from the target data distribution \(p_{\text{data}}(a^* | o)\). 
Traditional behavior cloning (BC) frames this as a single-step direct mapping \(a = f_\theta(o)\), 
optimizing the isotropic mean squared error \(\mathbb{E}_{o, a^* \sim p_{\text{data}}} \|f_\theta(o) - a^*\|_2^2\). 
While deployment is highly efficient, this objective reduces generation to an isolated endpoint prediction, 
failing to exploit the geometric structure of the expert action manifold.
Generative action policies address this by constructing a continuous trajectory over a virtual time \(t \in [0, 1]\), 
where \(a_t\) denotes the intermediate state. Starting from an initial source sample \(a_0 \sim p_0\), 
diffusion models progressively reach the final action via a denoising path, 
typically optimized using the denoising loss \(\mathcal{L}_{\text{DP}} = \mathbb{E}[\|\epsilon_\theta(a_t, o, t) - \epsilon\|_2^2]\). 
Evolving from this, flow matching directly formulates the trajectory via an ordinary differential equation (ODE) 
\(da_t = v_\theta(a_t, o, t) dt\) driven by a velocity field, 
optimizing \(\mathcal{L}_{\text{FM}} = \mathbb{E}[\|v_\theta(a_t, o, t) - \frac{d}{dt}a_t\|_2^2]\). 
This multi-step evolution naturally retains the capacity for progressive error correction.

To reduce inference latency, recent acceleration algorithms compress the intermediate path to directly learn a one-step endpoint mapping or a large-stride velocity. 
For instance, Consistency Models~\cite{prasad2024consistencypolicyacceleratedvisuomotor} learn a mapping \(f_\theta(o, a_t, t)\) by imposing a cross-time self-consistency constraint:
\(\mathcal{L}_{\text{CM}} = \mathbb{E}\left[ \left\| f_\theta(o, a_t, t) - f_{\phi}(o, a_{t'}, t') \right\|_2^2 \right]\),
enforcing any intermediate point on the trajectory to map to the same endpoint (where \(\phi\) is the target network parameter). 
Similarly, MeanFlow~\cite{geng2025meanflowsonestepgenerative} learns the average velocity \(u_\theta(a_t, o, r, t)\) over a time interval \([r, t]\) by satisfying an intrinsic identity between average and instantaneous velocities:
\(\mathcal{L}_{\text{MF}} = \mathbb{E}\left[ \left\| u_\theta(a_t, o, r, t) - \operatorname{sg}\left[ v_t - (t-r)\frac{d}{dt}u_\theta(a_t, o, r, t) \right] \right\|_2^2 \right]\),
where \(\operatorname{sg}[\cdot]\) denotes the stop-gradient operator.
However, whether relying on target-network distillation or higher-order derivative identities, the optimization processes encounter specific challenges. 
When the physical intermediate evolution is bypassed, 
reconstructing the corrective signal for the one-step mapping without relying on expensive cross-step constraints remains a crucial problem.

\paragraph{Explicit Drifting in Behavior Cloning.}
Drifting models offer a mathematical framework to isolate the correction mechanism entirely within the training phase. 
Given the target data distribution \(p_{\text{data}}\) and generated distribution \(q_\theta\), 
Drifting constructs an explicit continuous vector field \(V_{p,q}(a) = V_p^+(a) - V_q^-(a)\) comprising attraction and repulsion. 
This field dictates the optimal displacement for a generated sample \(a\) toward the data manifold. Based on a kernel function \(k(\cdot, \cdot)\), 
the attraction and repulsion fields are respectively defined as:
\begin{equation}
    V_p^+(a) = \frac{\mathbb{E}_{y\sim p_{\text{data}}}[k(a, y)(y - a)]}{\mathbb{E}_{y\sim p_{\text{data}}}[k(a, y)]}, \qquad V_q^-(a) = \frac{\mathbb{E}_{y\sim q_\theta}[k(a, y)(y - a)]}{\mathbb{E}_{y\sim q_\theta}[k(a, y)]}.
\end{equation}
The one-step generator \(f_\theta\) is trained to regress this field-displaced frozen target:
\begin{equation}
    \mathcal{L}_{\mathrm{Drift}}(\theta) = \mathbb{E}_{a_0 \sim p_0} \left\| f_\theta(o, a_0, 0) - \operatorname{sg}\left[ f_\theta(o, a_0, 0) + V_{p,q}(f_\theta(o, a_0, 0)) \right] \right\|_2^2.
\end{equation}
In unconditional generation, the rich sample sets supporting the distributions enable the vector field to provide smooth manifold guidance.
However, directly porting this explicit drifting mechanism to conditional behavior cloning encounters a fundamental structural degeneration. 
Consider the following proposition:

% \textbf{Proposition 1.} 
% \emph{
%     Assume that in behavior cloning, the dataset provides only a unique, 
%     deterministic expert action \(a_i^*\) for a specific observation \(o_i\), 
%     and the generator produces a single prediction \(a = f_\theta(o_i, a_0, 0)\) per training iteration. 
%     The empirical conditional distributions degenerate into Dirac measures \(\hat{p}_{\text{data}}(y | o_i) = \delta(y - a_i^*)\) 
%     and \(\hat{q}_\theta(y | o_i) = \delta(y - a)\). For any translation-invariant positive kernel \(k(\cdot, \cdot)\), 
%     the empirical conditional drifting field mathematically evaluates exactly to \(V_{\hat{p},\hat{q}}(a | o_i) = a_i^* - a\). 
%     Consequently, the field-displaced regression objective becomes exactly equivalent to the isotropic Mean Squared Error (MSE), 
%     losing all local geometric guidance.} 
%     \emph{(Proof. See Appendix A.)
% }
\begin{proposition}[Degeneration of Explicit Drifting in Behavior Cloning]
    \label{prop:degeneration_explicit_drifting}
    
    Assume that in behavior cloning, the dataset provides only a unique, 
    deterministic expert action \(a_i^*\) for a specific observation \(o_i\), 
    and the generator produces a single prediction \(a = f_\theta(o_i, a_0, 0)\) per training iteration. 
    The empirical conditional distributions degenerate into Dirac measures \(\hat{p}_{\text{data}}(y | o_i) = \delta(y - a_i^*)\) 
    and \(\hat{q}_\theta(y | o_i) = \delta(y - a)\). For any translation-invariant positive kernel \(k(\cdot, \cdot)\), 
    the empirical conditional drifting field mathematically evaluates exactly to \(V_{\hat{p},\hat{q}}(a | o_i) = a_i^* - a\). 
    Consequently, the field-displaced regression objective becomes exactly equivalent to the isotropic Mean Squared Error, 
    losing all local geometric guidance. (Proof. See Appendix~\ref{app:proof_prop1}.)
\end{proposition}

While the underlying control task may tolerate variations around \(a_i^*\) 
(e.g., alternative valid trajectories or recoverable states), 
the preceding analysis highlights that the empirical target distribution in standard BC training is too sparse to directly support the estimation of a continuous drifting field. 
To circumvent this degeneration in practice, models must resort to empirical approximations, 
such as synthesizing a localized empirical distribution by drawing multiple independent noise samples per observation during training and applying self-masking heuristics. 
While successful in avoiding the collapse to MSE, 
these approximations significantly increase computational overhead and entangle the correction field with transient sampling densities. 
This practical challenge motivates us to explore an alternative path: 
rather than regressing an explicit dynamic vector field, 
we extract the static local geometric structure directly from demonstrations.

\section{Implicit Drifting Policy}

Given the challenges of estimating explicit continuous vector fields under sparse empirical distributions, 
we explore an alternative path. Our core idea is to shift away from regressing dynamic drifting fields, 
and instead consider extracting the local geometric structure directly from static expert demonstrations, 
converting it into a potential objective that induces correction (see Fig~\ref{fig:pipeline}). 


\begin{figure}[t]
\centering

\includegraphics[width=\linewidth]{img/pipeline1.pdf}
% \caption{
% \textbf{Overview of Implicit Drifting Policy.}
% % IDP replaces explicit drifting-field regression with a geometry-aware potential induced by static demonstrations.
% (a) For each demonstration pair $(o_i,a_i^*)$, the expert action anchors the local action space.
% (b) Observation-conditioned weights $w_{ij}$ select expert actions under similar observations, forming a Conditional Expert Geometry $G_i$ around $a_i^*$.
% (c) Comparing $G_i$ with the reference geometry $\Sigma_{\mathrm{ref}}$ isolates the Local Geometry Excess $m_i$, visualized as the excess-constrained region and implemented as coordinate-wise precision excess.
% (d) The resulting metric $M_i$ defines a potential $E_i$, whose negative gradient $-\nabla E_i(a)$ supervises both the deployment proposal $y_i$ and the expert-proximal evaluation $z_i$ during training.
% }
\caption{
\textbf{Overview of Implicit Drifting Policy.}
(a) Expert action $a_i^*$ anchors the local action space.
(b) Weights $w_{ij}$ select observation-similar expert actions to form a Conditional Expert Geometry $G_i$ around $a_i^*$.
(c) Comparing $G_i$ with reference geometry $\Sigma_{\mathrm{ref}}$ yields the Local Geometry Excess $m_i$ (coordinate-wise precision excess).
(d) The induced metric $M_i$ defines a potential $E_i$, where $-\nabla E_i(a)$ supervises both deployment proposal $y_i$ and expert-proximal evaluation $z_i$.
}
\vspace{-2ex}
\label{fig:pipeline}
\end{figure}




\subsection{Implicit Correction Objective and Local-Aware Evaluation}
\label{sec:4.1}
To avoid regressing a high-dimensional and unstable continuous drifting field, 
we consider constructing an implicit correction field directed toward the expert target \(a_i^*\). 
In the most basic case, a simple attraction field can be written as \(V_{\mathrm{base}}(a) = a_i^* - a\). 
However, this simple attraction is isotropic. In multi-step generative models, 
isotropic targets, like MSE loss, are sufficient because the model can implicitly adjust and align with the action manifold through hundreds of inference iterations along the intermediate trajectory. 
However, for a one-step generator, it loses the ability to progressively correct errors along the intermediate path. 
If the training objective does not explicitly reinforce the tightness of constraints across different action dimensions, 
the prediction is highly prone to diverging from the valid manifold structure in a single stride.
Therefore, for a given predicted action \(a\), we consider applying a geometry-aware correction vector:
\begin{equation}
    \Delta_i(a) = M_i(a_i^* - a),
\end{equation}
where \(M_i \succeq 0\) is a positive semi-definite symmetric matrix acting as geometry-adaptive weights, 
determining the strength of the correction force along different directions. 
We elucidate the mathematical properties of this vector through the following proposition.

\begin{proposition}[Geometry-Induced Correction Potential]
    \label{prop:geometry_induced_potential}
    For any constant positive semi-definite symmetric matrix \(M_i \succeq 0\), the geometry-induced correction field defined by \(\Delta_i^{\text{geo}}(a) = M_i(a_i^* - a)\) is a conservative field. Mathematically, it corresponds exactly to the negative gradient (\(-\nabla_a E_i^{\text{geo}}(a) = \Delta_i^{\text{geo}}(a)\)) of the following geometry-only potential function:
    (Proof. See Appendix~\ref{app:proof_prop2}.)
    \begin{equation}
        E_i^{\text{geo}}(a) = \frac{1}{2} (a - a_i^*)^\top M_i (a - a_i^*).
    \end{equation}
\end{proposition}
The complete geometric potential objective incorporates the isotropic error of behavior cloning:
\begin{equation}
    E_i(a) = \frac{1}{2} \| a - a_i^* \|_2^2 + E_i^{\text{geo}}(a) = \frac{1}{2}[\underbrace{ \| a - a_i^* \|_2^2}_{\text{isotropic error}} + \underbrace{(a - a_i^*)^\top M_i (a - a_i^*)}_{\text{geometry-adaptive penalty}}]
\end{equation}
whose full negative gradient naturally yields a combined correction force: \(-\nabla_a E_i(a) = (I + M_i)(a_i^* - a)\) (Fig~\ref{fig:pipeline} (d)).
The first term of this potential is equivalent to the isotropic error of behavior cloning, 
while the second term adaptively increases the penalty for constrained directions based on \(M_i\). 
% Consequently, we transform the regression of an explicit vector field into the minimization of a potential.
We transform the regression of an explicit vector field into the minimization of a potential. 

However, the constraint characterized by matrix \(M_i\) is statistically grounded primarily within the local vicinity of the expert action \(a_i^*\). 
If the optimization process applies this potential constraint exclusively at the standard initial prediction point \(y_i = f_\theta(o_i, a_0, 0)\), 
the network might struggle to adequately learn the true topology of the local manifold. To ensure that this local potential effectively guides network learning, 
we consider introducing an expert-proximal evaluation anchor during training. 
By selecting a time scalar \(t_* \in (0, 1)\) close to \(1\), 
we construct a probing input \(\tilde{a}_i = a_i^* + (1 - t_*)\epsilon_i\) (where \(\epsilon_i \sim \mathcal{N}(0, I)\)) by injecting exploratory noise proportional to the distance around the expert action. 
We evaluate the mapping at this proximal location, yielding \(z_i = f_\theta(o_i, \tilde{a}_i, t_*)\). 
Therefore, the single-step prediction and the local probing are controlled by the same geometric potential objective, 
formulating our total training objective as:
% \begin{equation}
%     \mathcal{L} = \mathbb{E}_{i \sim \mathcal{D}} \left[ \tfrac{1}{t_*^2}\, E_i(y_i) + \tfrac{1}{(1-t_*)^2}\, E_i(z_i) \right],
% \end{equation}
% where the time-symmetric weights \(1/t_*^2\) and \(1/(1-t_*)^2\) are determined entirely by the schedule scalar \(t_*\). ...
\begin{equation}
    \mathcal{L} = \mathbb{E}_{i \sim \mathcal{D}} \left[ E_i(y_i) + \lambda_{\mathrm{prox}}\, E_i(z_i) \right],
    \label{eq:idp_total_loss}
\end{equation}
where \(\lambda_{\mathrm{prox}}\) is a hyperparameter balancing the standard endpoint regression and the expert-proximal evaluation. We adopt the implementation choice of MIP~\cite{pan2026adonoisingdispellingmyths} for this weight; details are deferred to Appendix~\ref{app:hyper}.
This construction establishes the skeleton of the entire implicit correction framework. 
The core challenge then becomes how to extract the adaptive constraint matrix \(M_i\) from static demonstration data to reflect the current control state.

\subsection{Conditional Expert Geometry as the Correction Source}

To ensure the adaptive penalty effectively limits invalid actions, 
we consider grounding the construction of matrix \(M_i\) in the authentic structure of the demonstration data. 
In control manifolds, similar observation states usually correspond to analogous robot poses or environmental layouts, 
which physically often impose similar operational bottlenecks and obstacle avoidance constraints. 
Therefore, under similar observation conditions, 
the variation in expert actions naturally reflects the specific motion constraints and acceptable tolerances of the current task.

To extract this structure, we need to identify and aggregate similar control contexts in the feature space. 
However, because feature representations dynamically update during training, 
to ensure that the similarity metric relies on stable features, we consider applying a stop-gradient operation and normalization to the embeddings before computation:
% \begin{equation}
%     h_i = \frac{\operatorname{sg}[\phi_\theta(o_i)]}{\|\operatorname{sg}[\phi_\theta(o_i)]\|_2 + \varepsilon}.
% \end{equation}
\begin{equation}
    h_i = \operatorname{sg}[\phi_\theta(o_i)]  \Big/ \,    \big(\|\operatorname{sg}[\phi_\theta(o_i)]\|_2 + \varepsilon \big).
\end{equation}
To evaluate the correlation between different states, 
we first compute the pairwise feature inner products \(s_{ij} = h_i^\top h_j\). 
Considering that the observation density across different regions of the state space in the demonstration dataset is often highly imbalanced
 (e.g., abundant samples in routine states but very few in edge cases), 
directly using inner products can lead to local weight collapse in sparse regions. 
To mitigate this issue, we adopt row-normalization \(\bar{s}_{ij} = \frac{s_{ij} - \mu_i}{\sigma_i + \varepsilon}\), 
and obtain the local neighborhood weights \(w_{ij} \propto \exp(\bar{s}_{ij})\) via Softmax.
Based on these weights, 
to quantify the degree of action divergence under the current observation, 
we define the Conditional Expert Geometry matrix within the local space of the expert action \(a_i^*\) (Fig~\ref{fig:pipeline} (b)):
\begin{equation}
    G_i = \mathbb{E}_{j \sim \mathcal{D} \setminus \{i\}} \left[ w_{ij} (a_j^* - a_i^*)(a_j^* - a_i^*)^\top \right].
    \label{eq:cond_expert_geometry}
\end{equation}
Note that the query sample itself (\(j=i\)) is explicitly excluded to prevent the zero-distance term from artificially collapsing the local variance estimate.
From a nonparametric statistical perspective, 
this construction can be interpreted through kernel smoothing, 
serving as an estimate of the local conditional covariance \(\Sigma_{a|o}\). 
Directions in \(G_i\) with small variance indicate strict operational constraints, 
whereas directions with large variance represent redundant dimensions of the action.

\subsection{Local Geometry Excess}

Directly utilizing the inverse of \(G_i\) as \(M_i\) could introduce biases. 
This is because certain action coordinates may naturally maintain extremely low variance globally across the entire task. 
We consider it necessary to distinguish the condition-specific constraints triggered purely by the current observation from this inherent global action scale.

In a general multivariate form, distinguishing the condition-specific constraints from the global marginal prior typically requires comparing their inverse covariance matrices. Let \(P_i^{\mathrm{cond}} = (G_i + \varepsilon I)^{-1}\) and \(P^{\mathrm{ref}} = (\Sigma_{\text{ref}} + \varepsilon I)^{-1}\). The geometry excess can be conceptualized via a non-negative truncated comparison, such as \(M_i \propto [(P^{\mathrm{ref}})^{-1/2} P_i^{\mathrm{cond}} (P^{\mathrm{ref}})^{-1/2} - I]_+\), which usually involves expensive generalized eigenvalue decomposition operations. Considering computational efficiency and to ensure relatively independent penalties across dimensions, we adopt a diagonal instantiation in practice. 
Specifically, we compute the diagonal conditional variance \(v_{i,d}^{\mathrm{cond}} = \mathbb{E}_{j \sim \mathcal{D}}[w_{ij} (a_{j,d}^* - a_{i,d}^*)^2]\) 
and the global reference variance \(v_d^{\mathrm{ref}} = \mathbb{E}_{i \sim \mathcal{D}} [(a_{i,d}^* - \bar{a}_d^*)^2]\). 
To eliminate dimensional impacts and compare the tightness of different dimensions on the same relative baseline, 
we convert them into normalized precision scales (i.e., standardized inverse variance):
% \begin{equation}
%     s_{i,d}^{\mathrm{cond}}
%     =
%     \frac{1/\sqrt{v_{i,d}^{\mathrm{cond}}+\varepsilon}}
%     {\mathbb{E}_{k}\big[1/\sqrt{v_{i,k}^{\mathrm{cond}}+\varepsilon}\big]},
%     \qquad
%     s_d^{\mathrm{ref}}
%     =
%     \frac{1/\sqrt{v_d^{\mathrm{ref}}+\varepsilon}}
%     {\mathbb{E}_{k}\big[1/\sqrt{v_k^{\mathrm{ref}}+\varepsilon}\big]}.
% \end{equation}
\begin{equation}
    s_{i,d}^{\mathrm{cond}} = \big( v_{i,d}^{\mathrm{cond}} + \varepsilon \big)^{-1/2} \Big/ \, \mathbb{E}_{k}\big[ (v_{i,k}^{\mathrm{cond}} + \varepsilon)^{-1/2} \big], 
    \quad
    s_d^{\mathrm{ref}} = \big( v_d^{\mathrm{ref}} + \varepsilon \big)^{-1/2} \Big/ \, \mathbb{E}_{k}\big[ (v_k^{\mathrm{ref}} + \varepsilon)^{-1/2} \big].
\end{equation}
Since we only wish to penalize dimensions that become more strictly constrained due to the current observation than they usually are, 
we consider utilizing the filtering property of the one-sided ReLU operator to define the Geometry Excess (Fig~\ref{fig:pipeline} (c)):
\begin{equation}
    m_{i,d}=\operatorname{ReLU}\left(s_{i,d}^{\mathrm{cond}} \Big/ \, (s_d^{\mathrm{ref}}+\varepsilon)-1\right),
    \qquad
    M_i=\operatorname{Diag}(m_i).
\end{equation}
The ReLU operator ensures that a dimension is extracted only when 
the local precision exceeds the global marginal precision prior. 
Consequently, we filter out inherent features that do not surpass the baseline, 
extract the specific geometry purely associated with the current state, 
and use it as the adaptive weight matrix \(M_i\) for the potential in Section~\ref{sec:4.1}.